% Budget: 2 pages, ~700 words, incl. Limitations and biases from the log
\chapter{Discussion}
\label{ch:discussion}

\section{Limitations and biases}
\label{sec:limitations}
% Moved from Positionality on 6 Nov, Celestin's own text.

This research is based on three main corpora, that introduce different perspectives, but also their own biases.
\begin{enumerate}
	\item First, the legal framework only takes into account the French one. The MIVILUDES provides a framework that is clear and documented, but also does not – of course – represent how these questions are tackled in other countries.
	\item Second, the interviews: Some interviews were conducted in French, and then translated to English for this paper. This methodology allowed me to gather rich and diverse data, but also introduced some limitations, and the responses were largely dependent on other factors: the context in which the interview was conducted, the mood of the interviewee and interviewer, and their willingness to share their thoughts and experiences. Some of the interviews were also conducted in the presence of other people, which could have influenced their responses.
	\item Third, the scholarly corpus is based on a systematic selection of articles and books, but according to criteria I thought were relevant to my research, and that I had to reshape as I went along, for example because scholars have long abandoned the term \textit{cult} in favor of other terms like \textit{new religious movements}, the criteria for relevant papers have been extended to other terms.
\end{enumerate}

Finally, absent from all of the data that was possible within the scope of this project is evidently the perspective of people that have been part of a group that could be conceived as a cult. This is a limitation of the research, and it is important to acknowledge that the perspectives of these people are not represented in this paper, and would represent a significant contribution to the understanding of the topic. However, it is also important to acknowledge that this perspective is not easily accessible, and that it would require a different methodology and a different set of ethical considerations.

Language bias across the three corpora
- Each epistemology comes from a narrow language representation, and this introduces its own bias.
- Legal framework (MIVILUDES):
    - Written in French.
    - Carries a specific cultural heritage, shaped by France's long-standing focus on laïcité.
- Scholarly corpus:
    - Mostly from the English-speaking world, or papers translated into English.
- Interviews:
    - Conducted in my immediate surroundings.
    - Participants tend to be English speakers based in Berlin, or French speakers based in France.
- The interaction between the corpora can also [unfinished: how?]
